\documentclass[11pt]{article}
\usepackage{amsmath,amssymb,amsthm,mathtools}
\usepackage{enumitem}
\usepackage[margin=1in]{geometry}
\newtheorem{theorem}{Theorem}
\newtheorem{lemma}{Lemma}
\newtheorem{definition}{Definition}
\newtheorem{proposition}{Proposition}
\newtheorem{warning}{Warning}
\newtheorem{corollary}{Corollary}
\newcommand{\K}{\mathsf K}
\newcommand{\A}{\mathsf A}
\newcommand{\E}{\mathsf E}
\newcommand{\Y}{Y^-}
\newcommand{\Wcmp}{W_{\mathrm{cmp}}}
\newcommand{\VZ}{V_Z}
\title{Step 76: Signed-to-Positive Weil Matching Gate}
\author{Six Birds / Formed-Layer Membrane Program}
\date{}
\begin{document}
\maketitle

\section{Purpose}
The RH membrane route requires a Douglas bridge
\[
  \A_Z\preceq \K_{\mathrm{cmp}},
  \qquad
  \A_Z=\VZ^*\VZ,\quad \K_{\mathrm{cmp}}=\Wcmp^*\Wcmp.
\]
Equivalently, there must be a contraction
\[
  \VZ=T\Wcmp,
  \qquad \|T\|\le 1.
\]
Step 75 proposed a log-side core and candidate positive feature slots
\[
  \Wcmp=(W_{\mathrm{pr}},W_{\infty},W_{\mathrm{pole}},W_{\mathrm{tail}})^T.
\]
The present note isolates the matching problem between the classical signed Weil explicit formula and the positive feature-map object required by the membrane framework.

\section{Signed forms and positive feature maps}

\begin{definition}[Signed Weil matching datum]
A signed matching datum on a real Hilbert response space \(\Y\) consists of Hermitian quadratic forms
\[
  q_Z[y]=\|\VZ y\|^2,
  \qquad
  q_{\rm EF}[y]=q_{\rm pr}^{\rm sgn}[y]+q_{\infty}^{\rm sgn}[y]+q_{\rm pole}^{\rm sgn}[y]+q_{\rm tail}^{\rm sgn}[y],
\]
and a proposed identity or inequality relating the zero-side anti-invariant readout to the completed explicit-formula side.
It is \emph{feature-positive} if there are Hilbert spaces \(\mathcal E_i\) and maps \(W_i:\Y\to\mathcal E_i\) such that
\[
  q_{\rm cmp}[y]
  =\sum_i \|W_i y\|^2
  \quad\text{and}\quad
  q_Z[y]\le q_{\rm cmp}[y]
  \quad\forall y\in\Y.
\]
\end{definition}

\begin{theorem}[Positive feature representation]
Let \(Q\) be a bounded Hermitian form on a Hilbert space \(\Y\). The following are equivalent:
\begin{enumerate}[label=(\roman*)]
\item \(Q\succeq0\).
\item There exists a Hilbert space \(\mathcal E\) and a bounded map \(W:\Y\to\mathcal E\) such that \(Q=W^*W\).
\end{enumerate}
In finite dimension one may take \(W=Q^{1/2}\).
\end{theorem}

\begin{proof}
If \(Q=W^*W\), then \(\langle y,Qy\rangle=\|Wy\|^2\ge0\). Conversely, if \(Q\succeq0\), the positive square root \(Q^{1/2}\) exists and gives \(Q=(Q^{1/2})^*Q^{1/2}\). The unbounded form version requires closed positive forms and a form-domain square-root representation.\qedhere
\end{proof}

\begin{corollary}[Signed slot obstruction]
A signed explicit-formula component cannot be interpreted as a positive feature slot unless its quadratic form is positive semidefinite on the declared domain. If a component is indefinite, it must either be grouped with other components into a positive completed block or carried as a defect.
\end{corollary}

\section{Douglas bridge as full-domain Weil positivity}

\begin{theorem}[Weil positivity equals Douglas domination after identification]
Suppose a completed Weil form on \(\Y\) is identified as
\[
  Q_W[y]
  =\langle y,(\K_{\rm cmp}-\A_Z)y\rangle,
  \qquad
  \A_Z=\VZ^*\VZ,
  \quad
  \K_{\rm cmp}=\Wcmp^*\Wcmp.
\]
Then the following are equivalent:
\begin{enumerate}[label=(\roman*)]
\item \(Q_W[y]\ge0\) for all \(y\in\Y\).
\item \(\A_Z\preceq\K_{\rm cmp}\).
\item There exists a contraction \(T\) with \(\VZ=T\Wcmp\).
\end{enumerate}
\end{theorem}

\begin{proof}
The equivalence of (i) and (ii) is just the definition of Loewner order. The equivalence of (ii) and (iii) is Douglas factorization: \(\VZ^*\VZ\preceq\Wcmp^*\Wcmp\) if and only if \(\VZ=T\Wcmp\) for some contraction \(T\).\qedhere
\end{proof}

\begin{warning}[What is not earned by the classical formula alone]
A signed explicit formula, a scalar trace equality, or positivity on a finite test family does not by itself produce \(\K_{\rm cmp}=\Wcmp^*\Wcmp\) with \(\A_Z\preceq\K_{\rm cmp}\). The matching record must identify the response space, the positive completed carrier features, the zero-side readout, and the domain/core on which the domination holds.
\end{warning}

\section{Minimal repair of a signed matching form}

Let
\[
  S=\K_{\rm vis}-\A_Z
\]
be the visible signed gap between a proposed carrier feature and the zero displacement. If \(S\nsucceq0\), there is a witness \(y\) with \(\langle y,Sy\rangle<0\), so the Douglas gate fails.

\begin{proposition}[Spectral positive repair, finite-dimensional version]
In finite dimension let \(S=U\operatorname{diag}(\lambda_i)U^*\). Define
\[
  S_-:=U\operatorname{diag}(\max\{-\lambda_i,0\})U^*,
  \qquad
  S_+:=U\operatorname{diag}(\max\{\lambda_i,0\})U^*.
\]
Then
\[
  S+S_-=S_+\succeq0.
\]
Thus a defect or completion feature of size at least the negative spectral part repairs the visible gap. This repair is a mathematical completion of the form, not automatically a lawful analytic-number-theory source record.
\end{proposition}

\begin{proof}
Immediate from the spectral decomposition.\qedhere
\end{proof}

\begin{warning}[Repair is not proof]
The formal repair \(S_-\) only identifies the obstruction. It is not an RH proof unless \(S_-\) is realized by upstream-visible prime, gamma, pole, tail, or strict-extension source records. Otherwise it is target-selected and therefore smuggled.
\end{warning}

\section{Core-to-full matching}

\begin{theorem}[Core matching gate]
Let \(\mathcal C\subset\Y_{\rm full}^-\) be a dense form core. Suppose forms
\[
  q_Z[y]=\|\VZ y\|^2,
  \qquad
  q_{\rm cmp}[y]=\|\Wcmp y\|^2
\]
are defined on \(\mathcal C\), with \(q_{\rm cmp}\) closable and \(\mathcal C\) a form core for its closure. If
\[
  q_Z[y]\le q_{\rm cmp}[y]\quad (y\in\mathcal C),
\]
and \(q_Z\) is continuous with respect to the closed \(q_{\rm cmp}\)-graph norm, then the inequality extends to the completed response space and yields a Douglas contraction.
\end{theorem}

\begin{proof}
For any element in the completion choose a core sequence converging in the closed graph norm. The inequality passes to the limit by closedness/continuity. The closed forms then satisfy \(\A_Z\preceq\K_{\rm cmp}\), hence Douglas factorization gives the contraction.\qedhere
\end{proof}

\section{Status taxonomy}

A proposed Weil matching gate receives one of the following statuses.
\begin{description}
\item[exact\_positive\_feature] Each required component is represented by a positive feature map and the Douglas domination holds on the completed response space.
\item[grouped\_positive\_feature] Individual signed components are not positive, but a completed grouping is positive and upstream-visible.
\item[defect\_paid] The missing positivity is represented by an explicit defect \(E_{\rm EF}\) that enters the obstruction budget.
\item[core\_only] Domination is proved on a dense core but closability/core-extension records are missing.
\item[finite\_support\_only] Domination is shown only on finite windows or finite tests.
\item[trace\_shadow] Only scalar trace or total accounting is controlled.
\item[failed\_matching] A negative witness for \(\K_{\rm cmp}-\A_Z\) is present.
\end{description}

\section{RH obligation after this step}
The concrete RH construction target is now:
\[
  \boxed{
  \text{construct an upstream-visible positive completed feature representation}
  }
\]
\[
  \Wcmp=(W_{\rm pr},W_{\infty},W_{\rm pole},W_{\rm tail})^T
\]
such that the completed Weil form is exactly
\[
  Q_W[y]=\|\Wcmp y\|^2-\|\VZ y\|^2
\]
on the full anti-invariant response space, or on a dense core with a full completion record. This is stronger than citing the signed explicit formula and weaker than RH only in the sense that the proof still also needs source-coercivity and fixed/exhaustive zero-ledger squeeze.

\end{document}
